\begin{scanpage}{0136}{118}
Consequently, an electric current filament located above an ideally conducting half-plane excites, at the point $P$, the field
\[E_z=ik^2p_z[u(d,\varphi'-\varphi'')-u(d,\varphi'+\varphi'')]\sqrt{\frac{2\pi}{kR}}e^{i(kR-\pi/4)},\tag{20.11}\]
and a magnetic current filament excites, at the point $P$, the field
\[H_z=ik^2m_z[u(d,\varphi'-\varphi'')+u(d,\varphi'+\varphi'')]\sqrt{\frac{2\pi}{kR}}e^{i(kR-\pi/4)}.\tag{20.12}\]
It is easy to see that the exponent $e^{ik[R-d\cos(\varphi'-\varphi'')]}$ in these expressions corresponds to the primary cylindrical wave arriving at the observation point $P$, and the exponent $e^{ik[R-d\cos(\varphi'+\varphi'')]}$ corresponds to the reflected cylindrical wave.

The moments $m_z$ and $p_z$ must be selected in such a way that in the direction $\phi''=\pi$ (Figure 46) the filament would create a field equal to the field of the primary edge wave above an infinite, ideally conducting plane. We will conclude these calculations in the following sections, but for now let us make still one other comment on the formulation of the problem.

In the previous chapters it was shown that the scattering object may be approximated by a series of sources --- ``luminous'' lines and points. Therefore, the problem of secondary diffraction may be formulated as a problem of searching for functions which describe the continuous change of the field of each such source during the passage through the boundary of the light and shadow corresponding to the source.
\sectiontitle{21}{Secondary Diffraction by a Strip (H-Polarization)}
A current filament with the moment $m_z$ which is positioned above an ideally conducting plane ($h=0$, Figure 46) creates in space the field
\end{scanpage}
